\chapter{Proofs for the Standard Small-Clause Transformation}
\label{annex:p12-small-clause-standard}

\begin{annexitem}{annex:p12.small.sat.closure}{thm:p12-small-sat-closure}{Coincidence and existential closure}
\begin{proof}
We first establish the valuation fact used here and in subsequent proofs.
The value of a term depends only on the variables occurring in it.
For a variable this is the definition of evaluation; a constant has no
valuation dependence. For \(f(t_1,\ldots,t_n)\), agreement on all its
variables gives agreement on each child's variables, hence equal child
values by induction and equal function values.

Formula truth depends only on its free variables. This follows for
atoms from the term result, and for logical constants without a
valuation condition. At each Boolean connective, agreement on the
union of the operands' free-variable sets gives equal operand truth
values and hence equal output truth values. For \(Qx B\), let
\(\nu,\mu\) agree on \(\operatorname{free}(B)\setminus\{x\}\).
For every \(d\in D\), the valuations \(\nu[x/d],\mu[x/d]\)
agree on \(\operatorname{free}(B)\): they agree at \(x\) by the
updates and elsewhere by the assumption. Induction gives the same
truth of \(B\) for each \(d\). Thus the set of witnesses is the
same for \(\exists\), and the all-values condition is the same for
\(\forall\). This proves coincidence for both quantifiers.

If \(\nu\) satisfies \(\Phi\), use its values
\(\nu(x_1),\ldots,\nu(x_k)\) as successive witnesses for the
existential closure. Conversely, truth of that finite existential
prefix supplies a tuple \((d_1,\ldots,d_k)\) under which its body
is true. Choose a fixed \(d_0\in D\), use the constant valuation
with value \(d_0\), and update it by the supplied tuple. Coincidence
ensures that this valuation satisfies \(\Phi\), regardless of
values assigned to other variables. No axiom of Choice is needed
for a finite string of witnesses. When \(k=0\), coincidence says
that truth of the sentence is independent of every valuation.

For the distinction between conventions, take \(D=\{0,1\}\),
\(P^{\mathfrak A}=\{0\}\), \(\nu(x)=0\) and \(\nu(y)=1\).
Then \(P(x)\land\neg P(y)\) is true. Its universal closure would
require \(P(d)\land\neg P(d)\) by setting both variables to any
one \(d\in D\), which is impossible.
\end{proof}
\end{annexitem}

\begin{annexitem}{annex:p12.small.context}{thm:p12-small-context}{Contextual replacement with bound variables}
\begin{proof}
Fix a structure in the specified model class. Order the two truth
values by \(0<1\), and write \(\llbracket A\rrbracket_\nu=1\)
when \(\mathfrak A,\nu\models A\), and \(0\) otherwise.
The assumption that \(A\) entails \(B\)
at every valuation is the pointwise inequality
\(\llbracket A\rrbracket_\nu\le
\llbracket B\rrbracket_\nu\). Conjunction and disjunction are
monotone in each operand. Negation reverses the order. Implication,
whose truth function is \(\neg a\lor b\), reverses the order in
its first operand and preserves it in the second.

Quantifiers are monotone operations on the valuation-dependent body.
If \(A\) is true at all \(\nu[x/d]\), the pointwise assumption
makes \(B\) true at all of them; this proves the universal case.
If \(A\) is true at one such update, the same update witnesses
\(B\); this proves the existential case. No assumption about the
selected variable's value being fixed outside the context is made.

Induct up the path from the selected occurrence to the root.
Each monotone step keeps the current implication direction, and
each antitone step reverses it. Their product is precisely the
specified nonzero polarity. For an equivalence assumption the
inserted truth values are equal at every valuation. Every Boolean
truth function, including biconditional, preserves equality of the
corresponding argument values, and both quantified truth conditions
preserve equality pointwise. Thus equality propagates even through
an ancestor biconditional and its zero-polarity descendants.
This proves replacement at all occurrences, including the root.

If equivalence is known only relative to background assumptions,
the argument applies to the structures satisfying those assumptions
at every valuation required by the replacement. Dropping them, or
checking the inserted formulas at just the current assignment before
entering a quantifier, would not meet the hypothesis.
\end{proof}
\end{annexitem}

\begin{annexitem}{annex:p12.small.simplify}{thm:p12-small-simplify}{All simplification rules and the normal-form invariant}
\begin{proof}
For a fixed valuation put \(a=\llbracket A\rrbracket\in\{0,1\}\).
Idempotence gives \(a\land a=a\) and \(a\lor a=a\).
Complementarity gives \(a\land\neg a=0\) and
\(a\lor\neg a=1\); reflexivity gives \(a\leftrightarrow a=1\)
and \(a\to a=1\). Negating truth or falsity interchanges them.
The conjunction identities with \(1,0\) yield \(a,0\), and the
disjunction identities yield \(1,a\), respectively. Using
\(a\to b=\neg a\lor b\), the four implication cases
\(a\to1,1\to a,a\to0,0\to a\) give \(1,a,\neg a,1\).
The two biconditional cases with \(1,0\) give \(a,\neg a\).
These are all eighteen Boolean schemes in the displayed collection.

For either quantifier, if \(x\notin\operatorname{free}(A)\),
coincidence from the preceding closure proof makes the truth value
of \(A\) the same under every \(\nu[x/d]\). Because \(D\)
is nonempty, both its universal and existential aggregation equal
that common value. This proves the two remaining schemes.
Contextual replacement proves equivalence at every application site.

Count every formula constructor and atomic formula as one tree node.
Writing \(N=N(A)\ge1\), the idempotence cases replace
\(2N+1\) nodes by \(N\); complementarity cases replace
\(2N+2\) by one; reflexivity cases replace \(2N+1\) by one.
Negated constants lose one node. A Boolean operation with a constant
has \(N+2\) nodes and is replaced by one, \(N\), or \(N+1\)
nodes. A redundant quantifier replaces \(N+1\) by \(N\).
Thus every actual rewrite strictly decreases a positive integer.
Commutative matching permutes operands only to recognize a rule;
it is not an additional rewrite that could be repeated forever.

Suppose an irreducible result has a logical constant at a proper
position. Its parent is a negation, binary connective, or quantifier.
In the first two cases an applicable constant scheme exists. In the
quantifier case its constant body has no free variable, so the
quantifier scheme applies. Each case contradicts irreducibility.
Therefore constants are absent or constitute the entire result.
Likewise an irreducible quantified subformula cannot have a vacuous
binder. No rule introduces a new free variable; rules such as
\(A\lor\neg A\mapsto\top\) can remove them. Finally, the
empty conjunction of clauses is true and a clause with no disjunct
is false, giving the stated outputs for \(\top\) and \(\bot\).
\end{proof}
\end{annexitem}

\begin{annexitem}{annex:p12.small.nnf}{thm:p12-small-nnf}{A decreasing interpretation for all NNF rewrites}
\begin{proof}
Both expansions of a biconditional have the same truth table as
that biconditional: they are true when both operands are true or
both false. De Morgan's laws, quantifier duality, the definition
of implication and double negation prove equivalence of all the
other rules. Quantifier duality follows directly from ``not every''
versus ``some not'', and ``not some'' versus ``every not'', for the
same domain values. Contextual replacement lifts these local facts.

Let the weights of \(A,B\) be \(a,b\), necessarily positive.
The two biconditional right sides have weights
\(5a+5b+5\) and \(3a+3b+3\), each strictly below
\(6a+6b+6\). For either De Morgan rule the left weight is
\(2a+2b+2\) and the right weight \(2a+2b+1\).
For either quantifier duality it is \(2a+2\) versus \(2a+1\).
Implication changes \(3a+2b+2\) to \(2a+b+1\), and double
negation changes \(4a\) to \(a\). These exhaust the rules.
Every constructor's weight function is strictly increasing in
each argument on positive integers. Consequently a decrease inside
any larger formula is a strict decrease of its total weight.
No infinite rewrite sequence is possible, including one that
duplicates arbitrarily complicated operands of a biconditional.

If a terminal formula contained a biconditional, choose one with
no biconditional ancestor. Its polarity is \(+1\) or \(-1\),
so one of the specified rules would apply. Implications are always
reducible. Once those two constructors are absent, a negation not
immediately on an atom is on another negation, a quantifier or a
conjunction/disjunction, and the corresponding rule applies.
Hence the terminal formula is NNF. Each rewrite retains precisely
the same names of free variables and of all occurring variables;
duplication changes occurrences but not these sets. Constant-free
input remains constant-free.

For the example, negative polarity gives
\[
\begin{aligned}
\neg(A\leftrightarrow B)
&\equiv \neg[(A\land B)\lor(\neg A\land\neg B)]\\
&\equiv \neg(A\land B)\land\neg(\neg A\land\neg B)\\
&\equiv(\neg A\lor\neg B)\land(\neg\neg A\lor\neg\neg B)\\
&\equiv(\neg A\lor\neg B)\land(A\lor B).
\end{aligned}
\]
If one instead expands every biconditional into two implications,
the corresponding sequence is
\[
\begin{aligned}
\neg(A\leftrightarrow B)
&\equiv\neg[(A\to B)\land(B\to A)]\\
&\equiv\neg[(\neg A\lor B)\land(\neg B\lor A)]\\
&\equiv\neg(\neg A\lor B)\lor\neg(\neg B\lor A)\\
&\equiv(A\land\neg B)\lor(B\land\neg A)\\
&\equiv(A\lor B)\land(A\lor\neg A)
       \land(\neg B\lor B)\land(\neg B\lor\neg A).
\end{aligned}
\]
For atomic operands these are four literal clauses, two tautological.
For compound operands, later NNF rewrites can obscure the original
syntactic complementary pair. The polarity-dependent expansion
avoids introducing these two particular redundant disjunctions.
\end{proof}
\end{annexitem}

\begin{annexitem}{annex:p12.small.miniscope}{thm:p12-small-miniscope}{Truth preservation and the quantifier-scope measure}
\begin{proof}
Fix a valuation and suppose \(x\) is not free in \(B\).
Its truth value is independent of the value assigned to \(x\).
If that value is true, then \(A\land B\) has the truth value
of \(A\), while \(A\lor B\) is true for every value of \(x\).
The displayed right sides have the same quantified truth values,
using nonemptiness for the existential quantification of a true
formula. If \(B\) is false, then \(A\lor B\) has the truth
value of \(A\), while \(A\land B\) is false for every value
of \(x\). Nonemptiness makes its universal quantification false.
This proves all four guarded schemes for both truth values of
\(B\). Exchanging operands is justified by commutativity.

For the two unguarded distributions, every value satisfies
\(A\land B\) exactly when every value satisfies \(A\) and
every value satisfies \(B\). Some value satisfies \(A\lor B\)
exactly when some value satisfies at least one of its operands.
These are the universal-conjunction and existential-disjunction
rules. Their additional free-variable conditions restrict when the
procedure uses these identities, not their truth. Contextual
replacement proves equivalence of each complete formula.

Write \(r=b(A)\), \(s=b(B)\). In a guarded step the selected
quantifier's contribution \(3^{r+s+1}\) is replaced by
\(3^r\); in a distributive step it is replaced by
\(3^r+3^s\). For \(r,s\ge0\),
\[
 3^r+3^s\le2\cdot3^{r+s}<3^{r+s+1}.
\]
Contributions of quantifiers inside \(A\) or \(B\) are unchanged.
An ancestor quantifier also has the same number of Boolean nodes
in its scope: the step moves or duplicates quantifiers but neither
duplicates nor removes the Boolean tree of \(A\circ B\).
Quantifiers outside that subtree are unaffected. Thus the total
nonnegative integer strictly decreases at every step. The sum is
zero in a quantifier-free formula, which has no available step.
This proves termination. It supplies no comparison between the
running times of different deduction procedures on the output.
\end{proof}
\end{annexitem}

\begin{annexitem}{annex:p12.small.renaming}{thm:p12-small-renaming}{The complete raw-renaming valuation lemma}
\begin{proof}
For a term variable \(z\), the hypothesis says precisely that
its two values agree after replacement by \(\rho(z)\).
Constants have unchanged interpretation. For a function application,
induction gives equal argument values, and applying the unchanged
function yields equal output values. This proves the term assertion.

For an atom, all term values agree, so its tuple belongs to the
same predicate interpretation on both sides; equality atoms obey
the same conclusion with the identity relation. Logical constants
are unchanged. For negation, the induction hypothesis gives equal
truth values before negating. For each of \(\land,\lor,\to\)
and \(\leftrightarrow\), agreement on the union of the operands'
free-variable sets yields equal operand values separately, hence
equal values under its specified Boolean truth function.

Consider a quantified formula \(Qx B\). For each \(d\in D\)
compare \(\nu_1[x/d]\) with \(\nu_2[\rho(x)/d]\).
At \(z=x\) their values agree modulo \(\rho\) by construction.
At \(z\in\operatorname{free}(B)\setminus\{x\}\), injectivity
on \(\operatorname{vars}(Qx B)\) gives
\(\rho(z)\ne\rho(x)\). Neither relevant coordinate changes,
and the original agreement on
\(\operatorname{free}(Qx B)\) applies. Therefore these updated
valuations agree modulo \(\rho\) on all free variables of \(B\).
Its induction hypothesis gives
\[
 \mathfrak A,\nu_1[x/d]\models B
 \quad\Longleftrightarrow\quad
 \mathfrak A,\nu_2[\rho(x)/d]\models\operatorname{raw}_\rho(B)
 \qquad(d\in D).
\]
For \(Q=\exists\), the same \(d\) is a witness in either
direction. For \(Q=\forall\), the equivalence holds for every
\(d\), and therefore the all-values conditions coincide.
This proves both quantified cases, including \(\rho(x)=x\).
Only injectivity on the occurring variables was used; no permutation
or injective identity-extension of the entire infinite alphabet is
required.
\end{proof}
\end{annexitem}

\begin{annexitem}{annex:p12.small.rectification}{thm:p12-small-rectification}{Fresh-block rectification, including repeated original binders}
\begin{proof}
Use induction on the number of formula-tree nodes, with the following
strengthened assertion for a call on \(A\) starting at \(n\).
Its unused names \(y_n,\ldots,y_{n+nq(A)-1}\) are pairwise
distinct and outside \(\operatorname{vars}(A)\). The output is
equivalent to \(A\), has the same free-variable set, has precisely
this block as its bound-name set with one name per binder, and
contains no other variables except free variables of \(A\).

An atom or logical constant has no binder and is returned unchanged.
At a negation, the child has fewer nodes, uses the same empty or
nonempty block, and its induction hypotheses pass through negation.
At a binary connective, the first child receives a block of length
\(nq(B)\) and the second the immediately following block of
length \(nq(C)\). The blocks are disjoint and fresh for both
original children. The induction hypotheses give equivalent
operands, their original free-variable sets, and the two disjoint
bound-name sets. None of these new names is free in the other
operand. The union of the blocks has the required total length,
and each of the four Boolean truth functions preserves equivalence.

At \(Qx B\), put \(y=y_n\) and
\(B'=\operatorname{raw}_{x\mapsto y}(B)\). This map is injective
on \(\operatorname{vars}(B)\), because \(y\) was absent.
It has the same tree shape and quantifier count as \(B\), and
its variables belong to \(\operatorname{vars}(B)\cup\{y\}\).
The later names in the supplied block are consequently fresh for
\(B'\). The recursive call on \(B'\) has fewer nodes than
\(Qx B\), so it terminates and satisfies the strengthened assertion.

For any valuation \(\nu\) and \(d\in D\),
\(\nu[x/d]\) and \(\nu[y/d]\) agree modulo the raw map on
\(\operatorname{free}(B)\). At \(x\) both values are \(d\);
all other relevant names are different from fresh \(y\) and are
unchanged. The preceding renaming lemma and then the recursive
equivalence give
\[
 \mathfrak A,\nu[x/d]\models B
 \ \Longleftrightarrow\ 
 \mathfrak A,\nu[y/d]\models B'
 \ \Longleftrightarrow\ 
 \mathfrak A,\nu[y/d]\models\operatorname{ren}(B',n+1).
\]
Taking an existential witness proves the \(\exists\) case, and
taking every \(d\) proves the \(\forall\) case.

Syntactic induction for raw renaming gives
\(\operatorname{free}(B')=
\rho(\operatorname{free}(B))\). At a quantifier this uses
injectivity to commute the image with removal of its bound name.
The recursive output has this same free-variable set, so binding
\(y\) removes exactly the image of \(x\) and leaves
\(\operatorname{free}(B)\setminus\{x\}\).
Its internal binders use exactly the later block, each once, while
the outer binder uses \(y_n\). The new bound names are absent
from the resulting free-variable set. All the strengthened claims
follow.

This argument permits the original input to reuse \(x\) at an
inner binder. The preliminary raw step then temporarily renames
that inner binder to \(y_n\) as well. The recursive call assigns
it a later fresh name and preserves its original binding behavior.
For example, \(\forall x(P(x)\lor\exists x Q(x))\) becomes
\(\forall y_1(P(y_1)\lor\exists y_2 Q(y_2))\), not a formula
that makes the two original binding occurrences identical.
\end{proof}
\end{annexitem}

\begin{annexitem}{annex:p12.small.skolem}{thm:p12-small-skolem}{Witness selection, substitution, expansion and reduct}
\begin{proof}
Let the selected subformula be \(E=\exists x B\) and its distinct
free parameters \(\bar z=(z_1,\ldots,z_m)\). For each tuple
\(\bar d\in D^m\), define
\[
 W_{\bar d}=\{e\in D:
       \mathfrak A,\nu[\bar z/\bar d][x/e]\models B\}.
\]
The arbitrary initial \(\nu\) does not affect this set: all free
variables of \(B\) are in \(\{x,z_1,\ldots,z_m\}\), and
coincidence applies. By Choice, choose \(e_{\bar d}\in W_{\bar d}\)
for each nonempty fiber. Fix one \(d_0\in D\), and interpret
the fresh symbol in an expansion \(\mathfrak A'\) by
\[
 f^{\mathfrak A'}(\bar d)=
 \begin{cases}
 e_{\bar d},&W_{\bar d}\ne\varnothing,\\
 d_0,&W_{\bar d}=\varnothing.
 \end{cases}
\]
For \(m=0\), \(D^0\) consists of the empty tuple and this is
the interpretation of a new constant. All old symbols keep their
interpretations, including any symbols introduced in earlier steps.

We spell out the substitution fact needed for this construction.
Put \(t=f(\bar z)\). Rectification ensures that no binder inside
\(B\) is named \(x\) or any \(z_i\). For an atom of \(B\),
term induction proves that substituting \(t\) for free \(x\)
and evaluating under \(\mu\) has the same effect as evaluating
the original terms under \(\mu[x/t^{\mathfrak A',\mu}]\).
The variable case is exactly this update; all other variables and
constants are unchanged; function applications follow from their
children. Predicate and equality truth then agree. Boolean
connectives preserve this agreement in all their arguments.
For an inner binder \(Qw C\), we have \(w\ne x\) and
\(w\notin\operatorname{vars}(t)\), so
\(t^{\mathfrak A',\mu[w/d]}=t^{\mathfrak A',\mu}\), and
updating \(w\) commutes with the displayed update of \(x\).
The induction hypothesis applies under every \(\mu[w/d]\).
Existential witnesses and universal all-values conditions therefore
coincide, proving
\[
 \mathfrak A',\mu\models B[x:=t]
 \quad\Longleftrightarrow\quad
 \mathfrak A',\mu[x/t^{\mathfrak A',\mu}]\models B.
\]

In the constructed expansion, if \(E\) is true under \(\mu\),
the corresponding fiber is nonempty and the chosen value of \(f\)
belongs to it. The substitution identity proves \(B[x:=t]\).
If the fiber is empty, no value satisfies \(B\), including the
default, so both \(E\) and \(B[x:=t]\) are false. Hence these
two formulas agree at every valuation in this expansion. The fresh
symbol does not occur in \(B\), and old-symbol truth agrees in
\(\mathfrak A\) and \(\mathfrak A'\). Contextual equivalence
therefore carries a model of the old whole sentence to a model of
the new one on the same domain.

Conversely, take any structure for the expanded signature. If
\(B[x:=t]\) is true, its evaluated term supplies an existential
witness by the same substitution identity. Thus
\(B[x:=t]\to\exists x B\) holds pointwise without Choice.
An existential occurrence in NNF has positive polarity: negations
occur only on atoms. Contextual monotonicity takes truth of the
new whole sentence to truth of the old sentence. Since the latter
contains no new symbol, it is true in the reduct.

An existential chosen with minimum position length has no
existential ancestor. Every binder above it is therefore universal,
and closedness plus rectification means that its free parameters
are precisely the occurring names bound by those ancestors.
Substitution cannot capture them at an inner binder. No quantifier
is duplicated; the selected existential is removed and the other
binder occurrences remain distinct. Closedness and rectification
are preserved. The number of existentials decreases by one,
which also proves termination of their successive elimination.

Freshness outside the current signature is what permits the
specified choice of \(f^{\mathfrak A'}\) without altering any old
interpretation. Merely being absent from the particular formula
does not permit such a change to a declared old symbol.
\end{proof}
\end{annexitem}

\begin{annexitem}{annex:p12.small.cnf}{thm:p12-small-cnf}{The two ordered phases of CNF formation}
\begin{proof}
All quantifiers now are universal. In a rectified closed formula,
if \((\forall x A)\circ B\) occurs, then \(x\) is not free
in \(B\). Otherwise that occurrence of \(x\), outside its unique
binder, would either be globally free or require a second binder
of the same name. Both contradict rectification and closedness.
The same argument works when the quantified operand is second.
The guarded identities are the reversals of the corresponding
miniscoping equivalences, so each pull preserves truth.

Count, for every quantifier occurrence, its Boolean ancestors and
sum those counts. Moving \(\forall x\) over its immediately
enclosing Boolean connective removes that one ancestor for the
selected binder. Descendant quantifiers inside \(A\) still have
the same Boolean connective as an ancestor, and quantifiers inside
\(B\) do too. No other Boolean ancestry changes. Thus the sum
decreases by exactly one. If some quantifier still lies beneath
a Boolean connective, there is a quantifier immediately below
a Boolean connective, possibly after selecting the outermost
quantifier in that branch. Its pull is allowed by the preceding
guard argument. Therefore a terminal result has all quantifiers
in one prefix and a quantifier-free matrix. No binder is duplicated,
so rectification is preserved throughout this phase.

For matrix distribution, every formula weight is at least two.
With operand weights \(a,b,c\), the rule's left weight is
\(a(b+c+1)=ab+ac+a\), and its right weight is
\(ab+ac+1\). Their difference \(a-1\) is positive.
Both constructor interpretations are strictly increasing on
positive arguments, so every contextual distribution strictly
decreases the total weight. The usual truth table gives
\(A\lor(B\land C)\) equivalent to
\((A\lor B)\land(A\lor C)\); commutative matching gives the
other operand order.

If a terminal matrix had a conjunction below a disjunction, choose
a lowest disjunction ancestor of a conjunction along such a path.
One can choose the last disjunction before the first conjunction;
they are parent and child, giving an available distribution.
Thus every disjunction has only literals below it. Flattening
the conjunctions and disjunctions yields a conjunction of clauses.
Universal quantification distributes across its finite conjuncts,
and variables absent from a conjunct may be removed vacuously.
This identifies the prefixed formula with the universally read
clause-set representation.

The scheduling qualification is real. Start with the rectified
formula \((\forall x P(x))\lor(Q\land R)\), where \(Q,R\)
are nullary predicates. Distributing prematurely and pulling in
each branch produces
\[
 [\forall x(P(x)\lor Q)]\land
 [\forall x(P(x)\lor R)].
\]
Pulling the first binder out gives
\(\forall x[(P(x)\lor Q)\land\forall x(P(x)\lor R)]\).
The remaining inner pull is blocked because \(x\) is free in
\(P(x)\lor Q\). Pulling the other binder first yields the
symmetric obstruction. Another rectification or a different
quantifier rule could repair this, but neither is part of the
displayed two-phase schedule. Pulling first avoids the duplication.
\end{proof}
\end{annexitem}

\begin{annexitem}{annex:p12.small.clause.simplify}{thm:p12-small-clause-simplify}{Clause instances, retained witnesses and condensation}
\begin{proof}
For an arbitrary term substitution \(\sigma\) and valuation
\(\nu\), define \(\nu_\sigma(x)=(x\sigma)^{\mathfrak A,\nu}\).
Term induction gives
\((t\sigma)^{\mathfrak A,\nu}=t^{\mathfrak A,\nu_\sigma}\):
the variable case is the definition; constants are unchanged;
and function applications use the induction hypotheses for their
arguments. Atom truth and its negation agree under these two
evaluations, and so does a disjunction of finitely many literals,
including the empty disjunction. Consequently a universally true
clause \(C\) has every instance \(C\sigma\) true under every
valuation.

If \(C\sigma\subseteq D\) as multisets, any true disjunct of
\(C\sigma\) is also present in \(D\). Thus universal truth
of \(C\) implies universal truth of \(D\). A model of the
clause set with \(D\) deleted still satisfies its retained witness
\(C\), and hence satisfies \(D\); the reverse implication is
immediate by removing a conjunct. A syntactic tautology is true
under every assignment because its atom and literal negation
have opposite truth values. Its deletion therefore preserves
the model class as well.

If \(C\) is obtained by removing duplicate literals from an
instance \(D\sigma\), the two disjunctions have the same truth
value: repetition adds no alternative truth value. Hence universal
truth of \(D\) implies that of \(C\). The additional required
subsumption of \(D\) by \(C\) gives the converse, by the
preceding paragraph. This proves condensation equivalence without
identifying multiset inclusion with ordinary set inclusion.

An injective renaming of a clause's occurring variables preserves
its universal truth: every valuation on either finite set of
variable names can be transferred to the corresponding names on
the other side, and the raw-renaming term/atom argument applies.
The remaining alphabet coordinates can be assigned any fixed
domain element. Different clauses may therefore be renamed apart
independently, because they are separate universally closed conjuncts.

Each clause deletion reduces the clause count by one and does not
increase the number of literal occurrences. Each proper condensation
strictly reduces the latter number and does not increase clause
count. Their sum is a nonnegative integer and strictly decreases
at every deletion or condensation. Renaming apart is a finite
representation step, not an endlessly repeatable reduction rule.

Finally, \(D=\{R(x,x),R(y,y)\}\) has two literal occurrences
under every substitution, even when their images coincide. It
cannot multiset-subsume a singleton. Substituting \(y\mapsto x\)
and deleting the duplicate gives \(C=\{R(x,x)\}\), which
subsumes \(D\) and is its equivalent condensation.
The retained-witness requirement excludes simultaneous deletion
of all clauses in a subsumption cycle; the next countermodels
exhibit why it is needed.
\end{proof}
\end{annexitem}

\begin{annexitem}{annex:p12.small.pipeline}{thm:p12-small-pipeline}{Composition of the standard transformation}
\begin{proof}
Existential closure is a finite operation with the fixed-structure
satisfaction equivalence proved above. Simplification terminates
and preserves equivalence; a constant result immediately has its
specified clause-set meaning. Otherwise its output is a finite
constant-free sentence. NNF rewriting and miniscoping each terminate
by their respective decreasing integer measures and preserve
equivalence. Indexed rectification terminates by tree-size induction
and leaves an equivalent closed NNF with unique bound names.

There are finitely many existential occurrences in that resulting
formula. Each successive local Skolem step removes exactly one
and introduces one fresh symbol of finite arity. Assuming Choice,
a model at a given step expands on the same domain to a model at
the next step. Repeating finitely often gives an expansion of the
original structure; every symbol of its original signature has
kept its interpretation. Conversely, the choice-free reduct
implication at each step composes backwards to a model of the
rectified sentence.

After those steps, prefix pulling and then matrix distribution
terminate, preserve equivalence in the expanded signature and
produce a finite universally quantified clause matrix. Clause
representation and the terminating retained-witness simplifications
preserve that expanded-signature model class. Thus every model of
the closed input has the claimed expansion, and every model of
the output has a reduct satisfying the closed input. Existential
closure then supplies precisely a satisfying valuation for the
original open input. For sentence input there are no free variables
to close and this is ordinary sentence satisfaction.

The full run is a finite concatenation of terminating finite phases.
Although intermediate formulas may grow, each individual rule has
a finite right side and each phase has the stated termination
argument. Only finitely many Skolem symbols are introduced, so
the extension of the possibly infinite original signature is finite.
The proof does not assert that arbitrary first-order satisfiability
has thereby been decided: that task is delegated to a deduction
procedure acting on the resulting clauses.
\end{proof}
\end{annexitem}

\begin{annexitem}{annex:p12.small.example}{thm:p12-small-example}{Every stage of the running example}
\begin{proof}
Write \(R_{xx}=R(x,x)\) only within this derivation. Associative
display of \(\land,\lor\) abbreviates the corresponding binary
trees; when needed, their reassociation is justified by their
truth functions rather than counted as a transformation step.
The input is already a constant-free NNF. The miniscoping steps are
\[
\begin{aligned}
&\forall x\exists y\forall z
       [R_{xx}\land((P(y)\lor R(x,y))\lor Q(z))]\\
&\equiv\forall x\exists y
       [R_{xx}\land\forall z((P(y)\lor R(x,y))\lor Q(z))]\\
&\equiv\forall x\exists y
       [R_{xx}\land((P(y)\lor R(x,y))\lor\forall z Q(z))]\\
&\equiv\forall x
       [R_{xx}\land\exists y((P(y)\lor R(x,y))\lor\forall z Q(z))]\\
&\equiv\forall x
       [R_{xx}\land((\exists y(P(y)\lor R(x,y)))\lor\forall z Q(z))]\\
&\equiv\forall x
       [R_{xx}\land((\exists y P(y)\lor\exists y R(x,y))\lor\forall z Q(z))]\\
&\equiv(\forall x R_{xx})\land
       \forall x[(\exists y P(y)\lor\exists y R(x,y))\lor\forall z Q(z)]\\
&\equiv(\forall x R_{xx})\land
       [\forall x(\exists y P(y)\lor\exists y R(x,y))\lor\forall z Q(z)]\\
&\equiv(\forall x R_{xx})\land
       [(\exists y P(y)\lor\forall x\exists y R(x,y))\lor\forall z Q(z)].
\end{aligned}
\]
The first two steps use that \(z\) is absent from \(R_{xx}\)
and from \(P(y)\lor R(x,y)\). The next uses that \(y\) is
absent from \(R_{xx}\); the following one uses its absence from
\(\forall z Q(z)\). Existential distribution is permitted because
\(y\) is free in both \(P(y)\) and \(R(x,y)\). Universal
distribution across the conjunction is permitted because \(x\)
is free in both conjuncts. The last two steps use its absence
from \(\forall z Q(z)\) and \(\exists y P(y)\), respectively.

Preorder allocation of five fresh binder names gives
\[
 (\forall y_1 R(y_1,y_1))\land
 [(\exists y_2 P(y_2)\lor
          \forall y_3\exists y_4 R(y_3,y_4))\lor\forall y_5 Q(y_5)].
\]
The first existential has no free parameter and is eliminated
using a fresh constant \(a\). The second has precisely the free
parameter \(y_3\), and is eliminated using a fresh unary \(f\).
The intermediate and resulting formulas are
\[
\begin{aligned}
 &(\forall y_1 R(y_1,y_1))\land
 [(P(a)\lor\forall y_3\exists y_4 R(y_3,y_4))\lor\forall y_5 Q(y_5)],\\
 &(\forall y_1 R(y_1,y_1))\land
 [(P(a)\lor\forall y_3 R(y_3,f(y_3)))\lor\forall y_5 Q(y_5)].
\end{aligned}
\]
These are expansion/reduct steps, not equivalences under arbitrary
preassigned interpretations of \(a,f\).

For prefix formation, first rewrite the bracketed right conjunct as
\[
\begin{aligned}
 &(\forall y_3[P(a)\lor R(y_3,f(y_3))])\lor\forall y_5Q(y_5)\\
 &\equiv\forall y_3[(P(a)\lor R(y_3,f(y_3)))\lor\forall y_5Q(y_5)]\\
 &\equiv\forall y_3\forall y_5[P(a)\lor R(y_3,f(y_3))\lor Q(y_5)].
\end{aligned}
\]
Pull \(y_1\) out of the left conjunct, then \(y_3\) and
\(y_5\) out of the right conjunct inside that prefix. Freshness
provides each side condition. The result is
\[
 \forall y_1\forall y_3\forall y_5
 [R(y_1,y_1)\land(P(a)\lor R(y_3,f(y_3))\lor Q(y_5))].
\]
Its matrix already is a conjunction of literal clauses. No
distribution is applicable to its outer conjunction. Removing
the displayed prefix in favor of universal clause reading and
renaming variables apart yields exactly
\(\{R(u,u)\}\) and \(\{P(a),R(v,f(v)),Q(w)\}\).

Without miniscoping, the existential in
\(\forall x\exists y\forall z
[R(x,x)\land(P(y)\lor R(x,y)\lor Q(z))]\)
has the sole free parameter \(x\). A single unary symbol \(g\)
gives the clauses
\(\{R(u,u)\}\) and
\(\{P(g(v)),R(v,g(v)),Q(w)\}\).
The miniscoped output has a constant in its \(P\)-literal but
two newly introduced symbols in total. This verifies the claimed
distinction between dependency arity and the number of symbols.
\end{proof}
\end{annexitem}

\begin{annexitem}{annex:p12.small.failures}{thm:p12-small-failures}{Finite countermodels for each forbidden transformation}
\begin{proof}
The following interpretations use the stated nonempty domains.
They specify all predicates occurring in their displayed examples;
other signature symbols can be interpreted arbitrarily.
\begin{enumerate}
\item For a missing free-variable guard, take \(D=\{0,1\}\)
and \(P^{\mathfrak A}=Q^{\mathfrak A}=\{1\}\).
Then \(\exists x(P(x)\land Q(x))\) is true, with witness \(1\).
The unguarded proposed result
\((\exists xP(x))\land Q(x)\) is false at \(\nu(x)=0\).
The moved-out \(Q(x)\) is no longer governed by the binder.

\item Take \(D=\{0,1\}\), \(P^{\mathfrak A}=\{0\}\)
and \(Q^{\mathfrak A}=\{1\}\). Every element satisfies
\(P(x)\lor Q(x)\), but neither predicate holds everywhere.
Thus \(\forall x(P(x)\lor Q(x))\) is true and
\((\forall xP(x))\lor(\forall xQ(x))\) is false.
In the same structure, neither element satisfies both predicates,
so \(\exists x(P(x)\land Q(x))\) is false, whereas
\((\exists xP(x))\land(\exists xQ(x))\) is true.
These refute the two missing unrestricted distributions separately.

\item For capture, let \(D=\{0,1\}\),
\(R^{\mathfrak A}=\{(0,1)\}\) and \(\nu(y)=1\).
The formula \(\exists xR(x,y)\) is true with witness \(0\).
Changing \(x\) to \(y\) while leaving the already occurring
\(y\) fixed yields \(\exists yR(y,y)\), false because the
relation has no diagonal pair. The proposed raw map is not
injective on the original occurring variables \(\{x,y\}\).

\item Let the declared signature already contain a constant \(a\),
even though \(a\) is absent from \(\exists xP(x)\). Take
\(D=\{0,1\}\), \(a^{\mathfrak A}=0\) and
\(P^{\mathfrak A}=\{1\}\). The existential is true, but
replacing its witness by the old \(a\) gives the false \(P(a)\).
An expansion must preserve \(a^{\mathfrak A}\), so cannot repair
this. A genuinely new constant could instead be interpreted by
\(1\), as the expansion theorem requires.

\item The distinct clauses \(\{P(x)\}\) and \(\{P(y)\}\)
subsume one another by variable renaming. On the one-element
domain \(D=\{0\}\) with \(P^{\mathfrak A}=\varnothing\),
their conjunction is false. Deleting both produces the empty
clause set, true in this structure. Sequential deletion with a
retained witness cannot perform that joint deletion.

\item For the running example, take \(D=\{0,1\}\), let
\(P^{\mathfrak A}=Q^{\mathfrak A}=D\), and let
\(R^{\mathfrak A}=\{(0,1),(1,0)\}\). Interpret \(a\) by
\(0\) and \(f\) by \(f(0)=1,f(1)=0\). Every diagonal
\(R(d,d)\) is false, while \(P(a)\), every \(Q(d)\) and
every \(R(d,f(d))\) are true. Hence all three pairwise
disjunction clauses in the statement are universally true.
The required unit \(\{R(u,u)\}\) is false, and the original
sentence is false as well. Adding the universal clause
\(\{\neg R(s,s)\}\) leaves the three erroneous clauses
satisfied in this same structure. In contrast, the original
sentence entails \(\forall sR(s,s)\), which is inconsistent
with \(\forall s\neg R(s,s)\) on every nonempty domain.
Thus the wrong replacement can destroy preservation of
satisfiability for the larger problem, not just fixed-model
equivalence.
\end{enumerate}
Nonempty-domain restrictions also have content. If empty domains
were allowed, \(\forall x\bot\) would be true while \(\bot\)
is false, and \(\exists x\top\) would be false while \(\top\)
is true. The vacuous-quantifier rules and the guarded identities
that use these cases would need corresponding qualifications.
\end{proof}
\end{annexitem}
